\documentclass[12pt,a4paper]{report}

\usepackage[utf8]{inputenc}
\usepackage[T1]{fontenc}
\usepackage[french]{babel}
\usepackage{graphicx}
\usepackage{geometry}
\usepackage{hyperref}
\usepackage{fancyhdr}
\usepackage{titlesec}

\geometry{margin=2.5cm}

\hypersetup{
    colorlinks=true,
    linkcolor=black,
    urlcolor=blue,
    citecolor=black,
}

% ------------------------------------------------------------
% PAGE DE GARDE
% ------------------------------------------------------------
\begin{document}

\begin{titlepage}
    \centering

   
    \includegraphics[width=5cm]{"C:/Users/ouatt/Desktop/logos/ensai_logo.png"}

    \vspace{0.5cm}
    {\Large \textsc{École nationale de la statistique}}\\
    {\Large \textsc{et de l'analyse de l'information}}

    \vspace{2cm}
    {\LARGE \textsc{Projet informatique}}\\[0.3cm]
    {\large Groupe 18}

    \vspace{1cm}
    {\LARGE \textsc{Rapport d'analyse et de conception}}

    \vspace{1.5cm}
    \rule{\linewidth}{0.5pt}

    \vspace{0.5cm}
  
    {\Huge \textbf{Concurrent SNCF}}

    \vspace{0.3cm}
    {\Large Une plateforme de vente de trajets ferroviaires concurrente de la SNCF}

    \vspace{0.5cm}
    \rule{\linewidth}{0.5pt}

    \vspace{2cm}
    \begin{flushleft}
        \textit{Étudiants :}\\
        YVANO Maxime\\
        VAN HECKE Jonas\\
        SIDOBRE Noé\\
        OUATTARA Youssouf Adam\\
        ARRAJI Salma
    \end{flushleft}

    \vfill

    \begin{flushright}
        \textit{Tuteur :}\\
        Kévin LEROY
    \end{flushright}

    \vspace{0.5cm}
    {Année 2026-2027}

\end{titlepage}

----------------------------------------------------------

\section*{Préambule}

Dans ce rapport, nous allons vous présenter \textbf{Concurrent SNCF}, notre
application de vente de trajets ferroviaires réalisée dans le cadre de notre
projet informatique en deuxième année. Celle-ci se base sur les données
fournies par l'API SNCF (Navitia) et a pour objectif de mettre en pratique
les principes de programmation et de modélisation acquis au cours de notre
scolarité.

\vspace{0.5cm}

\noindent\textbf{Contexte.} Nous sommes développeurs pour une toute nouvelle
société de transport ferroviaire qui exploite le réseau SNCF. Notre mission
est de mettre en place une solution permettant au service commercial de
notre entreprise de proposer des trajets en train à ses clients.

\vspace{0.3cm}

\noindent\textbf{Objectif.} Développer une API permettant, selon le niveau
d'habilitation de l'utilisateur (CLIENT, COLLABORATEUR ou ADMIN) :
\begin{itemize}
    \item la gestion des comptes utilisateurs et de leurs rôles ;
    \item la création et la gestion de lignes d'exploitation entre gares,
          existantes et récupérées via l'API SNCF ;
    \item la planification de trajets sur ces lignes, avec calcul
          automatique du temps de trajet via l'API SNCF ;
    \item la recherche de trajets par un client, avec aide à la saisie des
          gares par autocomplétion ;
    \item une fonctionnalité de différenciation commerciale par rapport au
          réseau SNCF existant.
\end{itemize}

\vspace{0.3cm}

\noindent Le détail de sa conception figure dans la suite de ce document.

\tableofcontents

% ------------------------------------------------------------
% CHAPITRE 1 : ORGANISATION
% ------------------------------------------------------------
\chapter{Organisation}

\section{Cahier des charges}

\subsection*{Présentation d'ensemble du projet Concurrent SNCF}



\subsection*{Utilisateurs de l'application}

Trois niveaux d'habilitation sont distingués dans l'accès aux fonctionnalités
de l'application : \textbf{CLIENT}, \textbf{COLLABORATEUR} et
\textbf{ADMIN}.



\subsection*{Fonctionnalités de base}


\begin{itemize}
    \item F1 --- Gestion des comptes utilisateur
    \item F2 --- Gestion des lignes d'exploitation
    \item F3 --- Gestion des trajets proposés
    \item F4 --- Recherche d'un trajet
    \item F5 --- Positionnement commercial (fonctionnalité différenciante)
\end{itemize}

\subsection*{Fonctionnalités optionnelles}

\begin{itemize}
    \item FO1 --- Réservation de trajet
    \item FO2 --- Classes de voyage
\end{itemize}

\subsection*{Outils de développement}



\section{Étapes de réalisation}



% ------------------------------------------------------------
% CHAPITRE 2 : MODÉLISATION UML
% ------------------------------------------------------------
\chapter{Modélisation UML}

\section{Diagramme de cas d'utilisation}

Ce diagramme de cas d'utilisation délimite la zone fonctionnelle du backend et structure les interactions selon le profil de l'utilisateur et avec l'API externe de la SNCF.
Les utilisateurs de l'application ont trois rôles, d'où le mécanisme d'héritage mis en place. Un "Utilisateur" se connecte ou créée son compte ce qui le renvoie vers son rôle, parmi les "Admins", les "Collaborateurs" et les "Clients", et ses possibilités dédiées. Le "Client" dispose des fonctionnalités de recherche, de consultation des tarifs et de réservation. Le "Collaborateur" s'occupe du réseau et de la planification des trajets. Et l'"Admin" supervise la gestion globale des rôles. Pour rechercher un trajet, agir sur les trajets et gérer les gares et lignes, il faut récupérer toutes les informations de l'API SNCF.

\section{Diagramme de base de données}


\section{Diagramme d'architecture}



\section{Diagramme de séquence}



\end{document}
